% ============================================================================
%  Appendix E -- Risposte alle domande (generated by tools/gen_appendices.py
%  from the slide decks, where they are the speaker notes of the quiz slides)
% ============================================================================
\chapter{Risposte alle domande}
\label{app:answers}

Ogni capitolo tematico e ogni capitolo tutorial termina con domande a scelta
multipla nello stile della prova intercorso (paragrafo \emph{Verificate di aver
capito}). Seguono le risposte: la lettera, il testo dell'opzione giusta, e una
riga di spiegazione. Prima rispondete, e poi leggete: una domanda di cui si è
letta la risposta non insegna niente.

\section*{Capitolo 1: L'installazione}
\begin{description}[leftmargin=3.2em,labelwidth=2.6em,itemsep=3pt]
\answer{1.1}{B}{\emph{ESP-IDF non è attivo in questo terminale.} L'attivazione dura quanto il terminale: eseguite export.sh, oppure aprite ESP-IDF PowerShell.}
\answer{1.2}{C}{\emph{PRG32-QT.} PRG32-QT interpreta la cartuccia e ha un debugger integrato; il compilatore serve solo per compilare le cartucce.}
\answer{1.3}{C}{\emph{chiama il runtime attraverso la stessa tabella di funzioni su ogni host.} Una cartuccia portabile dipende soltanto dall'ABI, la tabella di funzioni del runtime, che ogni host fornisce.}
\end{description}

\section*{Capitolo 2: Primo contatto}
\begin{description}[leftmargin=3.2em,labelwidth=2.6em,itemsep=3pt]
\answer{2.1}{B}{\emph{L'indirizzo del contatore.} la carica l'indirizzo dell'etichetta; lw legge poi in t1 il valore memorizzato a quell'indirizzo.}
\answer{2.2}{C}{\emph{ra, perché call lo sovrascrive.} call scrive in ra l'indirizzo di ritorno. Le chiamate della funzione chiamata lo sovrascriverebbero: una funzione non foglia salva ra sullo stack e lo ripristina prima di ret.}
\answer{2.3}{D}{\emph{44.} 300 = 256 + 44 = 1 0010 1100 in binario. AND con 1111 1111 conserva gli otto bit bassi: 0010 1100 = 44.}
\end{description}

\section*{Capitolo 3: Storia della computazione}
\begin{description}[leftmargin=3.2em,labelwidth=2.6em,itemsep=3pt]
\answer{3.1}{B}{\emph{Programma e dati stanno nella stessa memoria.} Il concetto di programma memorizzato: le istruzioni sono dati in memoria, prelevati ed eseguiti uno dopo l'altro.}
\answer{3.2}{C}{\emph{Il suo insieme di istruzioni è uno standard aperto, libero da realizzare.} RISC-V è un'architettura a standard aperto; chiunque può progettare e vendere un processore RISC-V senza licenza sull'architettura. Non è stata la prima RISC (RISC di Berkeley, MIPS e ARM sono precedenti).}
\answer{3.3}{C}{\emph{Ada Lovelace.} Ada Lovelace, per le sue note del 1843 sulla macchina analitica.}
\end{description}

\section*{Capitolo 4: Informazione digitale}
\begin{description}[leftmargin=3.2em,labelwidth=2.6em,itemsep=3pt]
\answer{4.1}{D}{\emph{256.} Un byte sono 8 bit: $2^{8}$ = 256 configurazioni, per esempio i naturali da 0 a 255.}
\answer{4.2}{B}{\emph{magenta.} 0xF81F = 11111 000000 11111: rosso e blu al massimo, verde a zero: magenta.}
\answer{4.3}{C}{\emph{5.} Quattro caratteri più il byte di terminazione a zero.}
\answer{4.4}{C}{\emph{0xCD.} RISC-V è little-endian: il byte meno significativo, 0xCD, va all'indirizzo più basso.}
\end{description}

\section*{Capitolo 5: Gli interi}
\begin{description}[leftmargin=3.2em,labelwidth=2.6em,itemsep=3pt]
\answer{5.1}{B}{\emph{-4.} Si complementa e si somma uno: 0000 0011 + 1 = 0000 0100 = 4, quindi la configurazione vale -4. (252 è la sua lettura senza segno.)}
\answer{5.2}{B}{\emph{100 + 28.} L'intervallo è [-128, 127]. 100 + 27 = 127 ci sta; 100 + 28 = 128 no: due positivi danno una configurazione negativa (1000 0000 = -128).}
\answer{5.3}{C}{\emph{t2 = 1, t3 = 0.} Con segno: -1 $<$ 1 è vero, quindi t2 = 1. Senza segno: 4\,294\,967\,295 $<$ 1 è falso, quindi t3 = 0.}
\answer{5.4}{A}{\emph{0,101.} 0,625 $\times$ 2 = 1,25 (1); 0,25 $\times$ 2 = 0,5 (0); 0,5 $\times$ 2 = 1,0 (1): 0,101 = 1/2 + 1/8.}
\end{description}

\section*{Capitolo 6: I reali}
\begin{description}[leftmargin=3.2em,labelwidth=2.6em,itemsep=3pt]
\answer{6.1}{D}{\emph{Eccesso k.} In eccesso k, X = x + k: a un numero maggiore corrisponde sempre una configurazione maggiore. Per questo codifica gli esponenti della virgola mobile, che vanno confrontati in fretta.}
\answer{6.2}{B}{\emph{1,5.} 0x180 = 384; 384 / 256 = 1,5. Parte intera 0x01, frazione 0x80 = 128/256 = 0,5.}
\answer{6.3}{B}{\emph{3,0.} 0100 0000 0100 0000 ...: S = 0, E = 1000 0000 = 128, F = 100...: $1,1_2\times 2^{128-127}=1,5\times 2=3,0$.}
\answer{6.4}{C}{\emph{infinito.} Tutti uni nell'esponente: infinito se la frazione è zero, NaN altrimenti.}
\end{description}

\section*{Capitolo 7: Verm1, parte 1}
\begin{description}[leftmargin=3.2em,labelwidth=2.6em,itemsep=3pt]
\answer{7.1}{B}{\emph{4.} 64/256 = 1/4 di cella per frame: una cella intera ogni 4 frame.}
\answer{7.2}{C}{\emph{(80, 32).} x = 10 $\times$ 8 = 80; y = 3 $\times$ 8 + 8 = 32.}
\answer{7.3}{B}{\emph{draw chiama altre funzioni, che sovrascrivono ra.} Una chiamata scrive in ra l'indirizzo di ritorno. Una funzione che fa chiamate deve mettere al sicuro sullo stack il proprio ra; a una funzione foglia non serve.}
\end{description}

\section*{Capitolo 8: Algebra di Boole: operatori e proprietà}
\begin{description}[leftmargin=3.2em,labelwidth=2.6em,itemsep=3pt]
\answer{8.1}{C}{\emph{x = 1, y = 1.} NAND è il complemento di AND: vale 0 solo quando entrambi gli ingressi valgono 1.}
\answer{8.2}{B}{\emph{$\nott{a}+b$.} $\nott{a\cdot\nott{b}}=\nott{a}+\nott{\nott{b}}=\nott{a}+b$, usando l'involuzione sul secondo termine.}
\answer{8.3}{C}{\emph{0x0000000A.} 0x5A = 0101 1010; AND con 0000 1111 conserva il nibble basso: 1010 = 0xA.}
\answer{8.4}{B}{\emph{ori t0, t0, 0x20.} OR con una maschera: x + 1 = 1 sul bit 5, x + 0 = x altrove. xori lo invertirebbe; addi propagherebbe un riporto nel bit 6 se il bit fosse già a 1.}
\end{description}

\section*{Capitolo 9: Algebra di Boole: funzioni e minimizzazione}
\begin{description}[leftmargin=3.2em,labelwidth=2.6em,itemsep=3pt]
\answer{9.1}{D}{\emph{256.} $2^{2^{3}}=2^{8}=256$.}
\answer{9.2}{B}{\emph{$A\,\nott{B}\,C$.} La riga 5 è 101: $A=1$, $B=0$, $C=1$, quindi il mintermine è $A\,\nott{B}\,C$.}
\answer{9.3}{B}{\emph{due variabili.} Un gruppo di $2^{k}$ celle elimina k variabili: quattro celle, due variabili.}
\answer{9.4}{C}{\emph{almeno due ingressi valgono 1.} $C_{out}$ è la funzione di maggioranza: $AB+C_{in}(A\oplus B)=AB+AC_{in}+BC_{in}$. Un numero dispari di 1 descrive invece il bit di somma S.}
\end{description}

\section*{Capitolo 10: Elettronica}
\begin{description}[leftmargin=3.2em,labelwidth=2.6em,itemsep=3pt]
\answer{10.1}{B}{\emph{260 $\Omega$.} R = (3,3 - 2) / 0,005 = 260 $\Omega$.}
\answer{10.2}{B}{\emph{il pMOS conduce e l'uscita vale 1.} Un gate basso chiude il pMOS, che collega l'uscita all'alimentazione; l'nMOS è aperto.}
\answer{10.3}{C}{\emph{6.} Un NAND (4) seguito da un invertitore (2).}
\answer{10.4}{B}{\emph{Perché il piedino legga un 1 definito quando il pulsante è rilasciato.} Senza, il piedino rilasciato sarebbe flottante; il pull-up lo lega a 3,3 V, e l'interruttore premuto prevale portandolo a massa.}
\end{description}

\section*{Capitolo 11: Macchine combinatorie}
\begin{description}[leftmargin=3.2em,labelwidth=2.6em,itemsep=3pt]
\answer{11.1}{B}{\emph{5.} $\lceil\log_2 26\rceil$ = 5, perché $2^{4}$ = 16 $<$ 26 $\le$ 32 = $2^{5}$. Restano sei parole non assegnate: il codice è incompleto.}
\answer{11.2}{C}{\emph{8 uscite, una a 1.} $2^{3}$ = 8 uscite, e ne è attiva esattamente una: la rappresentazione decodificata (one-hot) della parola in ingresso.}
\answer{11.3}{B}{\emph{3.} n linee di selezione scelgono tra $2^{n}$ ingressi: 3 linee per 8 ingressi.}
\answer{11.4}{C}{\emph{non rileva l'errore.} Due alterazioni lasciano invariata la parità. Un solo bit di parità rileva soltanto un numero dispari di errori.}
\end{description}

\section*{Capitolo 12: Verm1, parte 2}
\begin{description}[leftmargin=3.2em,labelwidth=2.6em,itemsep=3pt]
\answer{12.1}{C}{\emph{0x0a.} 0x08 OR 0x02 = 0x0a: un bit per pulsante.}
\answer{12.2}{B}{\emph{DOWN (1).} 3 XOR 1 = 2: DOWN è l'opposto di UP.}
\answer{12.3}{B}{\emph{ogni elemento della tabella è di 4 byte.} Gli elementi sono parole: l'elemento k sta allo scostamento 4k.}
\answer{12.4}{C}{\emph{t1 e t4 sono uguali.} XOR vale 1 dove i bit differiscono: nessuna differenza, tutti zeri.}
\end{description}

\section*{Capitolo 13: Macchine sequenziali}
\begin{description}[leftmargin=3.2em,labelwidth=2.6em,itemsep=3pt]
\answer{13.1}{B}{\emph{dal solo stato presente.} Moore: w va da S a U. In una macchina di Mealy l'uscita dipende dallo stato e dall'ingresso.}
\answer{13.2}{C}{\emph{Q resta a 1.} S = R = 0 è la condizione di mantenimento: la retroazione conserva il valore memorizzato.}
\answer{13.3}{C}{\emph{Il flip-flop cambia solo su un fronte del clock.} Il latch è trasparente per tutto il livello attivo del clock; il flip-flop campiona D solo sul fronte.}
\answer{13.4}{B}{\emph{4.} $\lceil\log_2 10\rceil$ = 4 variabili di stato, con il codice binario di lunghezza minima della Capitolo~11. Ne servirebbero dieci con un codice one-hot.}
\end{description}

\section*{Capitolo 14: Registri e bus}
\begin{description}[leftmargin=3.2em,labelwidth=2.6em,itemsep=3pt]
\answer{14.1}{B}{\emph{Perché solo la sorgente abilitata piloti le linee.} Più sorgenti condividono le stesse linee; tutte tranne una devono essere elettricamente scollegate (alta impedenza), altrimenti due uscite si contenderebbero il filo.}
\answer{14.2}{C}{\emph{32.} Un bit per impulso: n impulsi per n bit. Un trasferimento parallelo richiede un impulso e n fili.}
\answer{14.3}{C}{\emph{8192.} $2^{10}$ = 1024 parole di 8 bit: 8192 bit, cioè 1 KiB.}
\end{description}

\section*{Capitolo 15: La memoria}
\begin{description}[leftmargin=3.2em,labelwidth=2.6em,itemsep=3pt]
\answer{15.1}{C}{\emph{64 KiB.} $2^{16}$ = 65\,536 locazioni di un byte: 64 KiB, lo spazio di indirizzi di un Commodore 64.}
\answer{15.2}{D}{\emph{0x01.} Little-endian: il byte meno significativo, 0x01, sta all'indirizzo più basso.}
\answer{15.3}{A}{\emph{Perché i suoi condensatori perdono la carica.} Ogni bit è una carica in un condensatore che si disperde; un circuito interno alla memoria riscrive periodicamente ogni cella.}
\answer{15.4}{B}{\emph{Con sw all'indirizzo di un registro del dispositivo.} RISC-V non ha istruzioni di I/O: i dispositivi sono mappati in memoria e si raggiungono con normali load e store.}
\end{description}

\section*{Capitolo 16: La memoria permanente}
\begin{description}[leftmargin=3.2em,labelwidth=2.6em,itemsep=3pt]
\answer{16.1}{D}{\emph{Flash.} La flash è l'evoluzione a blocchi della EEPROM. Alla EPROM serve la luce ultravioletta; ROM e PROM non si possono cancellare.}
\answer{16.2}{B}{\emph{Il dato viene prelevato dalla memoria centrale e copiato nella cache, con i suoi vicini.} L'accesso è semplicemente più lento: il blocco che contiene l'indirizzo viene caricato nella cache, scommettendo che presto serviranno indirizzi vicini.}
\answer{16.3}{B}{\emph{La latenza.} La latenza è l'attesa che il settore arrivi sotto la testina; il tempo di posizionamento dipende dal movimento del braccio.}
\answer{16.4}{C}{\emph{una transizione tra pit e land.} Un cambio di riflettività, in un verso o nell'altro, è un 1; nessun cambio è uno 0.}
\end{description}

\section*{Capitolo 17: Verm1, parte 3}
\begin{description}[leftmargin=3.2em,labelwidth=2.6em,itemsep=3pt]
\answer{17.1}{C}{\emph{254, 255, 0, 1.} Dalla coda, head\_i - length + 1 = -2, cioè 254 dopo AND 255, fino alla testa: 254, 255, 0, 1.}
\answer{17.2}{B}{\emph{il byte 127.} 3 $\times$ 40 + 7 = 127.}
\answer{17.3}{B}{\emph{la partita comincia una volta.} Solo il frame in cui A scende dà s1 = 0x10: un fronte, una transizione.}
\answer{17.4}{B}{\emph{perché la testa possa entrare nella cella lasciata dalla coda.} Nello stesso passo la coda si sposta: quella cella è libera per la testa.}
\end{description}

\section*{Capitolo 18: Il processore}
\begin{description}[leftmargin=3.2em,labelwidth=2.6em,itemsep=3pt]
\answer{18.1}{B}{\emph{il registro di stato con i flag.} RISC-V non ha codici di condizione: i salti confrontano direttamente i registri.}
\answer{18.2}{C}{\emph{solo load e store.} È un'architettura load-store: lw, lh, lb, sw, sh, sb e le forme senza segno.}
\answer{18.3}{B}{\emph{t0 è stato inserito sullo stack.} Il puntatore scende di una parola e il valore è memorizzato sulla nuova cima: un push.}
\answer{18.4}{B}{\emph{accumulatore := accumulatore + [y].} L'accumulatore è l'operando implicito e la destinazione.}
\end{description}

\section*{Capitolo 19: Istruzioni e interruzioni}
\begin{description}[leftmargin=3.2em,labelwidth=2.6em,itemsep=3pt]
\answer{19.1}{B}{\emph{i sette meno significativi.} I bit da 6 a 0, in ogni formato: il processore sa sempre dove guardare per prima cosa.}
\answer{19.2}{C}{\emph{4.} Gli indirizzi contano i byte e un'istruzione è lunga quattro byte.}
\answer{19.3}{C}{\emph{mepc.} mepc, scritto dall'hardware alla trap e usato da mret. mtvec contiene l'indirizzo del gestore.}
\answer{19.4}{C}{\emph{viene prelevata un'istruzione illegale.} È causata dall'istruzione in esecuzione, in modo sincrono; le altre vengono dall'esterno, in un momento qualunque.}
\end{description}

\section*{Capitolo 20: I bus e la scheda video}
\begin{description}[leftmargin=3.2em,labelwidth=2.6em,itemsep=3pt]
\answer{20.1}{C}{\emph{921\,600 byte.} 307\,200 pixel $\times$ 3 byte = 921\,600 byte.}
\answer{20.2}{B}{\emph{100 MB/s.} 2 byte $\times$ 50 milioni di trasferimenti al secondo = 100 MB/s.}
\answer{20.3}{C}{\emph{la posizione del suo colore in una palette.} Il byte è un indice; il colore vero si cerca in una tabella di 256 voci, la palette.}
\answer{20.4}{B}{\emph{L'uso dei processori grafici per calcoli non grafici.} General-purpose computing on graphics processing units: calcolo generico su unità di elaborazione grafica.}
\end{description}

\section*{Capitolo 21: Ingresso e uscita}
\begin{description}[leftmargin=3.2em,labelwidth=2.6em,itemsep=3pt]
\answer{21.1}{B}{\emph{Trasforma una combinazione di riga e colonna in un codice di tasto.} Da una posizione decodificata nella matrice dei tasti a un codice compatto di n bit, come il codificatore della Capitolo~11.}
\answer{21.2}{B}{\emph{tre celle con filtri rosso, verde e blu.} Tre sub-pixel, ciascuno una cella di cristalli liquidi dietro un filtro colorato; un TFT in ognuno ne conserva lo stato.}
\answer{21.3}{B}{\emph{I pixel sono pilotati uno alla volta e nel frattempo devono mantenere il proprio stato.} Persistenza: senza un transistor per pixel, le celle rispondono lentamente e conservano male il proprio stato mentre si scandiscono le altre.}
\answer{21.4}{C}{\emph{polling.} Il programma interroga periodicamente il dispositivo. Con le interruzioni sarebbe il dispositivo a chiamare il programma.}
\end{description}

\section*{Capitolo 22: Verm1, parte 4}
\begin{description}[leftmargin=3.2em,labelwidth=2.6em,itemsep=3pt]
\answer{22.1}{B}{\emph{32 30 35.} I caratteri '2', '0', '5': 0x32, 0x30, 0x35. Le unità vanno allo scostamento 2.}
\answer{22.2}{B}{\emph{8.} Un passo ogni 4 frame è 64/256: 32 + 4 $\times$ score = 64, quindi score = 8.}
\answer{22.3}{B}{\emph{5.} la sono due istruzioni, auipc e addi; poi lw, beqz e ret.}
\answer{22.4}{B}{\emph{non chiama funzioni e usa solo registri a e t.} Una funzione foglia che tocca solo registri caller-saved non ha nulla da salvare.}
\end{description}

\section*{Capitolo 23: Il Cartridge Store}
\begin{description}[leftmargin=3.2em,labelwidth=2.6em,itemsep=3pt]
\answer{23.1}{B}{\emph{pubblicate la 1.0.1.} Una versione pubblicata si considera permanente: una correzione è una nuova versione con lo stesso id.}
\answer{23.2}{B}{\emph{i metadati vanno dentro una cartuccia, che ha una sola architettura.} L'elenco dei file appartiene al bundle; una singola cartuccia dichiara la propria architettura, data sulla riga di comando.}
\answer{23.3}{C}{\emph{in una variabile d'ambiente della vostra shell.} Un token è una credenziale: non va mai in un file condiviso o sotto controllo di versione.}
\end{description}

\section*{Capitolo 24: Oltre un solo core}
\begin{description}[leftmargin=3.2em,labelwidth=2.6em,itemsep=3pt]
\answer{24.1}{B}{\emph{5.} Con n illimitato, S = 1 / (1 - p) = 1 / 0,2 = 5.}
\answer{24.2}{C}{\emph{le false dipendenze tra usi dello stesso nome di registro.} Due valori indipendenti che condividono un nome ricevono registri fisici diversi e si possono calcolare nello stesso momento.}
\answer{24.3}{B}{\emph{30.} 256 bit sono 32 byte per giro: 960 / 32 = 30. Il codice non cambia.}
\answer{24.4}{B}{\emph{OR.} L'uscita vale 1 ogni volta che almeno un ingresso vale 1: OR.}
\end{description}
